\documentclass[11pt]{article}
\usepackage[margin=1in]{geometry}
\usepackage{amsmath, amssymb}
\usepackage{braket}
\usepackage{booktabs}
\usepackage{enumitem}
\usepackage[compact]{titlesec}
\usepackage[hidelinks]{hyperref}

\setlist[itemize]{left=1.25em, itemsep=2pt}
\setlist[enumerate]{left=1.25em, itemsep=2pt}

\begin{document}
\begin{center}
  {\Large When Do Variational Quantum Circuits\\ Keep Federated Training Data Private?}\\[6pt]
  CSCI-739 Quantum Machine Learning, Final Project Proposal\\[2pt]
  AJ Barea, \texttt{ajb6289@rit.edu}, September 24, 2026
\end{center}

\section*{Motivation}

In federated learning, clients train a shared model without sending their data: each client
computes a gradient on its own examples and sends only that gradient to a server, which
aggregates them~\cite{mcmahan2017}. On classical neural networks this protection is weak. An
attacker holding the model and one client's gradient can optimize a dummy input until its
gradient matches the observed one, recovering the private training example~\cite{zhu2019}.

Whether the same attack works on variational quantum circuits (VQCs) is an open question, and
the literature disagrees. Kumar et al.~\cite{kumar2023} argue that expressive data encodings
make inversion intractable: recovering the input means solving systems of high-degree
multivariate Chebyshev polynomials, and an underparameterized attacker faces exponentially many
spurious local minima. Papadopoulos et al.~\cite{papadopoulos2025} then reconstruct inputs from
VQC gradients numerically, including batch-trained data, once the model is sufficiently
overparameterized. Heredge et al.~\cite{heredge2025} connect privacy to the circuit's dynamical
Lie algebra and find that the properties that make a VQC trainable also make its input easier to
extract. This project measures where each of these claims holds.

\section*{Research questions}

\begin{enumerate}[label=\textbf{RQ\arabic*.}, left=2.5em]
  \item At a fixed parameter count, how does attack success change between a more expressive
        encoding (more data re-uploading layers~\cite{perezsalinas2020}) and a deeper trainable
        block?
  \item How do batch size, an unknown label, and a trained (rather than random) model change
        attack success?
  \item Does realistic device noise degrade the attack enough to act as a defense, and what
        does it cost in trainability?
\end{enumerate}

\section*{Method}

\begin{itemize}
  \item \textbf{Model.} Angle encoding with $R_Y$ rotations, repeated $r$ times (data
        re-uploading), interleaved with $\ell$ trainable blocks of $R_Y R_Z$ rotations and a CNOT
        chain. The output is the mean of $\braket{Z}$ over all qubits, so every input feature
        reaches the gradient. Four qubits first, then eight.
  \item \textbf{Threat model.} An honest-but-curious server knows the circuit and its parameters
        and observes one client's gradient.
  \item \textbf{Attack.} Gradient matching~\cite{zhu2019}: minimize the squared distance between
        the observed gradient and the gradient of a candidate input, using L-BFGS-B with random
        restarts, differentiating through the simulator.
  \item \textbf{Outcomes.} Each attack is classified as \emph{recovered} (input found, distance
        measured modulo $2\pi$), \emph{ambiguous} (a different input with the same gradient), or
        \emph{stuck} (a local minimum). I report success rates with ranges over many random
        seeds, and count only parameters whose gradient is nonzero.
  \item \textbf{Noise.} Shot noise, then Qiskit Aer noise models built from the calibration data
        of IBM Heron devices (\texttt{ibm\_kingston}, \texttt{ibm\_fez}, \texttt{ibm\_marrakesh}).
  \item \textbf{Tools.} PennyLane for the attack, since it needs gradients of gradients; Qiskit
        Aer for device noise. Code:
        \url{https://github.com/ajbarea/csci739-qfl-privacy}.
\end{itemize}

\section*{Preliminary result}

A first run on four qubits (one example, known label, untrained model, five seeds each)
already separates the three outcomes:

\begin{center}
\begin{tabular}{@{}lcccc@{}}
  \toprule
  Encoding reps $\times$ trainable layers & Parameters & Recovered & Ambiguous & Stuck \\
  \midrule
  $1 \times 1$ & 8   & 1 (+1 within 0.07) & 2 & 1 \\
  $1 \times 4$ & 32  & 5 & 0 & 0 \\
  $4 \times 1$ & 32  & 3 & 0 & 2 \\
  $4 \times 4$ & 128 & 4 & 0 & 1 \\
  \bottomrule
\end{tabular}
\end{center}

At equal parameter count, the more expressive encoding resisted the attack more often than the
deeper trainable block, consistent with~\cite{kumar2023}. Five seeds are too few to conclude
anything; the full study will establish whether this effect holds.

\section*{Deliverables}

A LaTeX report documenting the method and findings, centered on a privacy map of attack success
over encoding expressivity and parameter count, and an in-class presentation in weeks 12--14.

\begin{thebibliography}{9}
\setlength{\itemsep}{2pt}

\bibitem{mcmahan2017}
B.~McMahan, E.~Moore, D.~Ramage, S.~Hampson, and B.~Ag\"uera y Arcas.
Communication-efficient learning of deep networks from decentralized data.
In \emph{Proc.\ AISTATS}, PMLR 54:1273--1282, 2017.

\bibitem{zhu2019}
L.~Zhu, Z.~Liu, and S.~Han.
Deep leakage from gradients.
In \emph{Advances in Neural Information Processing Systems 32 (NeurIPS)}, 2019.

\bibitem{kumar2023}
N.~Kumar, J.~Heredge, C.~Li, S.~Eloul, S.~H.~Sureshbabu, and M.~Pistoia.
Expressive variational quantum circuits provide inherent privacy in federated learning.
arXiv:2309.13002, 2023.

\bibitem{papadopoulos2025}
G.~Papadopoulos, S.~Eloul, Y.~Satsangi, J.~Heredge, N.~Kumar, C.-F.~Chen, and M.~Pistoia.
A numerical gradient inversion attack in variational quantum neural-networks.
arXiv:2504.12806, 2025.

\bibitem{heredge2025}
J.~Heredge, N.~Kumar, D.~Herman, S.~Chakrabarti, R.~Yalovetzky, S.~H.~Sureshbabu, C.~Li, and
M.~Pistoia.
Characterizing privacy in quantum machine learning.
\emph{npj Quantum Information}, 11:80, 2025.

\bibitem{perezsalinas2020}
A.~P\'erez-Salinas, A.~Cervera-Lierta, E.~Gil-Fuster, and J.~I.~Latorre.
Data re-uploading for a universal quantum classifier.
\emph{Quantum}, 4:226, 2020.

\end{thebibliography}

\end{document}
